\begin{figure}[htbp]
\centering
\begin{tikzpicture}
    % Convex function (left plot)
    \begin{scope}[xshift=-4.5cm]
        % Axes
        \draw[->] (-2,0) -- (2.5,0) node[right] {$x$};
        \draw[->] (0,-0.5) -- (0,3.5) node[above] {$f(x)$};

        % Convex function (parabola)
        \draw[thick,UofSCGarnet,domain=-1.8:1.8,smooth,samples=50] plot (\x,{0.5*\x*\x + 0.5});

        % Two points on curve
        \coordinate (A) at (-1.2,{0.5*1.44 + 0.5});
        \coordinate (B) at (1.4,{0.5*1.96 + 0.5});
        \fill[UofSCAtlantic] (A) circle (2pt);
        \fill[UofSCAtlantic] (B) circle (2pt);

        % Chord connecting points
        \draw[thick,UofSCAtlantic,dashed] (A) -- (B);

        % Midpoint on chord
        \coordinate (M) at (0.1,{0.5*(0.5*1.44 + 0.5 + 0.5*1.96 + 0.5)});
        \fill[UofSCHorseshoe] (M) circle (2pt);

        % Midpoint on curve
        \coordinate (Mc) at (0.1,{0.5*0.01 + 0.5});
        \fill[UofSCHorseshoe] (Mc) circle (2pt);

        % Arrow showing chord above curve
        \draw[->,UofSCHorseshoe,thick] (Mc) -- ($(Mc)!0.8!(M)$);

        % Global minimum marker
        \fill[UofSCGarnet] (0,0.5) circle (2.5pt);
        \node[below,UofSCGarnet,font=\small] at (0,0.3) {global min};

        % Title
        \node[above,font=\bfseries] at (0,3.3) {Convex};
    \end{scope}

    % Non-convex function (right plot)
    \begin{scope}[xshift=4.5cm]
        % Axes
        \draw[->] (-2.5,0) -- (2.5,0) node[right] {$x$};
        \draw[->] (0,-0.5) -- (0,3.5) node[above] {$f(x)$};

        % Non-convex function
        \draw[thick,UofSCGarnet,domain=-2:2,smooth,samples=100] plot (\x,{0.3*\x*\x*\x*\x - 1.2*\x*\x + 2});

        % Two points on curve
        \coordinate (A2) at (-1.5,{0.3*5.0625 - 1.2*2.25 + 2});
        \coordinate (B2) at (1.5,{0.3*5.0625 - 1.2*2.25 + 2});
        \fill[UofSCAtlantic] (A2) circle (2pt);
        \fill[UofSCAtlantic] (B2) circle (2pt);

        % Chord connecting points
        \draw[thick,UofSCAtlantic,dashed] (A2) -- (B2);

        % Show chord goes below curve at origin
        \coordinate (Mchord) at (0,{0.5*(0.3*5.0625 - 1.2*2.25 + 2 + 0.3*5.0625 - 1.2*2.25 + 2)});
        \coordinate (Mcurve) at (0,2);
        \fill[UofSCRose] (Mcurve) circle (2pt);

        % Arrow showing chord below curve
        \draw[->,UofSCRose,thick] (Mcurve) -- ($(Mcurve)!0.8!(Mchord)$);

        % Local minima markers
        \fill[UofSCHorseshoe] (-1.41,{0.3*4 - 1.2*2 + 2}) circle (2.5pt);
        \fill[UofSCHorseshoe] (1.41,{0.3*4 - 1.2*2 + 2}) circle (2.5pt);
        \node[below left,UofSCHorseshoe,font=\scriptsize] at (-1.41,{0.3*4 - 1.2*2 + 2 - 0.2}) {local min};
        \node[below right,UofSCHorseshoe,font=\scriptsize] at (1.41,{0.3*4 - 1.2*2 + 2 - 0.2}) {local min};

        % Local maximum marker
        \fill[UofSCRose] (0,2) circle (2.5pt);
        \node[above,UofSCRose,font=\scriptsize] at (0,2.2) {local max};

        % Title
        \node[above,font=\bfseries] at (0,3.3) {Non-convex};
    \end{scope}
\end{tikzpicture}
\caption{Comparison of convex and non-convex functions. \textbf{Left:} A convex function (parabola) where the chord connecting any two points lies above the curve, and the unique critical point is the global minimum. \textbf{Right:} A non-convex function with multiple local minima and a local maximum. The chord connecting the two blue points passes below the curve (violating convexity). Gradient descent may converge to any local minimum depending on initialization.}
\label{fig:convex_nonconvex}
\end{figure}
